\documentclass[10pt,a4paper]{amsart}
\usepackage[margin=19mm]{geometry}
\usepackage{amsmath,amssymb,mathtools}
\usepackage[scr=rsfs,cal=euler]{mathalfa}
\usepackage{xfrac}
\usepackage[unicode,hidelinks,bookmarks=false]{hyperref}
\newcommand{\ie}{i.e.\ }
\newcommand{\cf}{cf.\ }
\newcommand{\resp}{respectively\ }
\newcommand{\Subsection}{\subsection}
\providecommand{\nobreakdash}{\nobreak\hbox{-}}
\setlength{\emergencystretch}{2em}
\multlinegap0pt
\usepackage{amssymb}
\usepackage[shortlabels]{enumitem}
\newtheorem{lemma}{Lemma}
\newtheorem{theorem}{Theorem}
\title{A data-driven model reduction approach for backward fractional diffusion-wave equations}
\author{Dakang Cen, Zhiyuan Li, Wenlong Zhang}
\date{Source: arXiv:2604.22211v1 (CC BY 4.0)}
\begin{document}
\maketitle

\noindent\emph{Source and licence.} A data-driven model reduction approach for backward fractional diffusion-wave equations. Version arXiv:2604.22211v1 (CC BY 4.0), distributed by arXiv under the Creative Commons Attribution 4.0 International licence, http://creativecommons.org/licenses/by/4.0/, as recorded in the arXiv licence metadata for this exact version and shown on the abstract page https://arxiv.org/abs/2604.22211v1. Full licence notice: \texttt{licenses/arxiv-2604.22211-CC-BY-4.0.txt}.\par\smallskip

\noindent\textit{Editorial scope.} Counted unit: the article's theorem error-adj (source0.tex lines 397--402). The article's complete original proof of it is delivered with it (source0.tex lines 403--469, one \texttt{proof} environment of 67 lines). No mathematics is written, repaired or re-derived: the statement and the proof are the article's own lines, in the article's own notation, and no retained line is substituted or re-flowed.  Retained uncounted as the source-local context the proof needs: lines 354--359; lines 362--367; lines 371--376; lines 378--386; lines 388--394.  Nothing was omitted from any retained span, so the package carries no omission record.  The retained lines cite 2 works, whose entries are transcribed verbatim from the article's own bibliography source in the e-print and delivered with short editorial labels recorded in the item's reference mapping.  No retained line contains a figure environment, an image-inclusion command or drawing code, so no image is delivered and the package declares no asset dependency.  Editorial formatting only, in addition: an AMS \texttt{amsart} preamble that restates the article's own declarations and macros used by the retained lines, namely the environments \texttt{lemma}, \texttt{theorem} on separate counters, together with the standard packages those lines need.  The retained environments print in this document as Lemma 1, Theorem 1 in order of appearance, whereas the article prints the same items with the numbering of its own full document.  The complete native source of the article as distributed in the arXiv e-print for this exact version is bundled unchanged as \texttt{sources/199078/source0.tex}.
\begin{equation}\label{num-scheme-fed}
\begin{cases}
  (\bar\partial_t^\nu \widetilde{V}_h^n,\varphi) + (\nabla \widetilde{U}_h^n,\nabla \varphi) = (q(x)\omega_{2-\alpha}(t_n),\varphi), & \forall \varphi\in X_h, \\
  (\nabla\widetilde{V}_h^n,\nabla\varphi)  = (\bar\partial_t^\nu \nabla\widetilde{U}_h^n,\nabla\varphi),
\end{cases}
\end{equation}

\begin{equation}\label{num-scheme-pod}
\begin{cases}
  (\bar\partial_t^\nu \widetilde{V}_m^n,\varphi) + (\nabla \widetilde{U}_m^n,\nabla \varphi) = (q(x)\omega_{2-\alpha}(t_n),\varphi), & \forall \varphi\in X_h^m, \\
  (\nabla\widetilde{V}_m^n,\nabla\varphi)  = (\bar\partial_t^\nu \nabla\widetilde{U}_m^n,\nabla\varphi).
\end{cases}
\end{equation}

\begin{lemma}\label{FD-ineq}
(\cite{LiaoH2018L1})For $V^n$, $1\leq n\leq N$, one has
\begin{align*}
(\bar\partial_t^\nu V^n,V^n)\geq \frac{1}{2}\bar\partial_t^\nu \|V^n\|_{L^2}^2.    
\end{align*}
\end{lemma}

\begin{lemma}\label{grownwall}
(\cite{LyuP2022SFOR})Let $(g^n)_{n=1}^N$ and $(\lambda_l)_{l=0}^{N-1}$ be given nonnegative sequences. Assume that there exists a constant $\Lambda$ such that $\Lambda\geq\sum_{l=0}^{N-1}\lambda_l$, and that the maximum step satisfies 
$$\max_{1\leq n \leq N}\tau_n\leq \frac{1}{\sqrt[\nu]{\Gamma{(2-\nu)}\Lambda}}.$$
Then, for any nonnegative sequence $(v^k)_{k=0}^N$ and $(w^k)_{k=0}^N$ satisfying
$$\sum_{k=1}^nA_{n-k}^{(n)}\nabla_\tau\big[(v^k)^2+(w^k)^2\big]\leq\sum_{k=1}^n\lambda_{n-k}\big(v^k+w^k\big)^2+(v^n+w^n)g^n,~~1\leq n\leq N,$$
it holds that
$$v^n+w^n\leq 4E_\nu(4\Lambda t_n^\nu)\bigg(v^0+w^0+\max_{1\leq k\leq n}\sum_{j=1}^kP_{k-j}^{(k)}g^j\bigg),~~1\leq n\leq N,$$
where $E_\nu(z)=\sum_{k=0}^\infty\frac{z^k}{\Gamma{(1+k\nu)}}$ is the Mittag-Leffler function.
\end{lemma}

\begin{lemma}\label{P}
For the sequence $(P_{n-j}^{(n)})_{j=1}^n$, some properties are given in \cite{LiaoH2018L1}.
\begin{align*}
&\sum_{j=k}^nP_{n-j}^{(n)}A_{j-k}^{(j)}\equiv1,~~1\leq k\leq n,\\
&0\leq P_{n-j}^{(n)}\leq \Gamma{(2-\nu)}\tau_j^\nu,~~\sum_{j=1}^nP_{n-j}^{(n)}t_j^{-\alpha/2}\leq C,~~ 1\leq j\leq n\leq N.  
\end{align*}
\end{lemma}

\begin{theorem}\label{error-adj}
Let $\widetilde{U}_h^n$ and $\widetilde{U}_m^n$ be the solutions of (\ref{num-scheme-fed}) and (\ref{num-scheme-pod}). Then there holds 
\begin{align*} 
\frac{1}{N}\sum_{n=1}^N\|\widetilde{U}_h^n-\widetilde{U}_m^n\|_{L^2}^2\leq CN^{2-\alpha}\sum_{j=m+1}^r\lambda_j.    
\end{align*} 
\end{theorem}

\begin{proof}
Introducing the Ritz projection operator $R_n^m:X_h\rightarrow X_h^m$, such that $(\nabla R_n^m \widetilde{U}_h^n, \nabla \varphi_m)=(\nabla \widetilde{U}_h^n, \nabla \varphi_m)$, $\forall \varphi_m\in X_h^m\subset X_h$.  Splitting $e_m^n=(\widetilde{U}_h^n-R_h^m\widetilde{U}_h^n)+(R_h^m\widetilde{U}_h^n-\widetilde{U}_m^n):=\rho_u^n+\theta_u^n$, $z_m^n=(\widetilde{V}_h^n-R_h^m\widetilde{V}_h^n)+(R_h^m\widetilde{V}_h^n-\widetilde{V}_m^n):=\rho_v^n+\theta_v^n$. 
First, we give an estimate for $\rho_u^n$.
\begin{align}\label{pod-1}
\frac{1}{N}\sum_{n=1}^N\|\rho_u^n\|_{L^2}^2\leq \frac{1}{N}\sum_{n=1}^N\|\rho_u^n\|_{H_0^1}^2\leq C\sum_{j=m+1}^r\lambda_j.    
\end{align}
Next, we derive an estimate for $\theta_u^n$.
\begin{align*}
&(\bar\partial_t^\nu\theta_v^n,\varphi_m)+(\nabla\theta_u^n,\nabla\varphi_m)\\
=&(\bar\partial_t^\nu R_h^m\widetilde{V}_h^n,\varphi_m)+(\nabla R_h^m\widetilde{U}_h^n,\nabla\varphi_m)-(\bar\partial_t^\nu \widetilde{V}_m^n,\varphi_m)-(\nabla \widetilde{U}_m^n,\nabla\varphi_m)\\
=&(\bar\partial_t^\nu R_h^m\widetilde{V}_h^n,\varphi_m)+(\nabla \widetilde{U}_h^n,\nabla\varphi_m)-(q(x)\omega_{2-\alpha}(t_n),\varphi_m)\\
=&(\bar\partial_t^\nu (R_h^m\widetilde{V}_h^n-\widetilde{V}_h^n),\varphi_m)\\
:=&-(\bar\partial_t^\nu \rho_v^n,\varphi_m).
\end{align*}
Substituting $\varphi_m=\theta_v^n$ into the above equation, it gives that
\begin{align}\label{pod-2}
&(\bar\partial_t^\nu\theta_v^n,\theta_v^n)+(\nabla\theta_u^n,\nabla\theta_v^n)=-(\bar\partial_t^\nu \rho_v^n,\theta_v^n).
\end{align}
Furthermore, one has
\begin{align*}
(\nabla \theta_v^n,\nabla\varphi_m)
&=(\nabla R_h^m\widetilde{V}_h^n,\nabla\varphi_m)-(\nabla \widetilde{V}_m^n,\nabla\varphi_m)\\
&:=(\nabla \widetilde{V}_h^n,\nabla\varphi_m)-(\bar\partial_t^\nu \nabla\widetilde{U}_m^n,\nabla\varphi_m)\\
&=(\bar\partial_t^\nu \nabla(\widetilde{U}_h^n-\widetilde{U}_m^n),\nabla\varphi_m)\\
&=(\bar\partial_t^\nu \nabla (\rho_u^n+\theta_u^n),\nabla\varphi_m).
\end{align*}
Substituting $\varphi_m=\theta_u^n$ into the above equation, we have
\begin{align}\label{pod-3}
(\nabla \theta_v^n,\nabla\theta_u^n) 
=(\bar\partial_t^\nu \nabla \rho_u^n,\nabla\theta_u^n)+(\bar\partial_t^\nu \nabla \theta_u^n,\nabla\theta_u^n).
\end{align}
For (\ref{pod-3}), from (\ref{pod-2}), we get
\begin{align}\label{pod-4}
(\bar\partial_t^\nu\theta_v^n,\theta_v^n)+(\bar\partial_t^\nu \nabla \theta_u^n,\nabla\theta_u^n)=-(\bar\partial_t^\nu \nabla \rho_u^n,\nabla\theta_u^n)-(\bar\partial_t^\nu \rho_v^n,\theta_v^n).
\end{align}
From Lemma \ref{FD-ineq}, for (\ref{pod-4}), we have
\begin{align*}
\frac12\bar\partial_t^\nu\|\theta_v^n\|_{L^2}^2+\frac12\bar\partial_t^\nu \|\nabla \theta_u^n\|_{L^2}^2
&\leq \|\bar\partial_t^\nu \nabla \rho_u^n\|_{L^2}\|\nabla\theta_u^n\|_{L^2}+\|\bar\partial_t^\nu \rho_v^n\|_{L^2}\|\theta_v^n\|_{L^2}\\
&\leq (\|\bar\partial_t^\nu \nabla \rho_u^n\|_{L^2}+\|\bar\partial_t^\nu \rho_v^n\|_{L^2})(\|\nabla\theta_u^n\|_{L^2}+\|\theta_v^n\|_{L^2}).
\end{align*}
Lemma \ref{grownwall} gives that 
\begin{align*}
\|\theta_u^n\|_{L^2}\leq C\big(\|\theta_v^n\|_{L^2}+ \|\nabla \theta_u^n\|_{L^2}\big) 
&\leq C\max_{1\leq k\leq n}\sum_{j=1}^kP_{k-j}^{(k)}(\|\bar\partial_t^\nu \nabla \rho_u^j\|_{L^2}+\|\bar\partial_t^\nu \rho_v^j\|_{L^2}),~1\leq n\leq N.
\end{align*}
Suppose $\max_{1\leq k\leq n}\sum_{j=1}^kP_{k-j}^{(k)}(\|\bar\partial_t^\nu \nabla \rho_u^j\|_{L^2}+\|\bar\partial_t^\nu \rho_v^j\|_{L^2})=\sum_{j=1}^{n_0}P_{n_0-j}^{({n_0})}(\|\bar\partial_t^\nu \nabla \rho_u^j\|_{L^2}+\|\bar\partial_t^\nu \rho_v^j\|_{L^2})$, $n_0\leq n$, from Lemma \ref{P}, one has
\begin{align*}
\|\theta_u^n\|_{L^2}^2
&\leq C\bigg(\sum_{j=1}^{n_0}P_{n_0-j}^{({n_0})}(\|\bar\partial_t^\nu \nabla \rho_u^j\|_{L^2}+\|\bar\partial_t^\nu \rho_v^j\|_{L^2})\bigg)^2\\
&\leq C\bigg(\sum_{j=1}^{n_0}\tau_j^{\alpha/2}(\|\bar\partial_t^\nu \nabla \rho_u^j\|_{L^2}+\|\bar\partial_t^\nu \rho_v^j\|_{L^2})\bigg)^2\\
&\leq C\bigg(\sum_{j=1}^{n}\tau_j^{\alpha/2}(\|\bar\partial_t^\nu \nabla \rho_u^j\|_{L^2}+\|\bar\partial_t^\nu \rho_v^j\|_{L^2})\bigg)^2\\
&\leq C\sum_{j=1}^{n}\tau_j^{\alpha}\sum_{j=1}^{n}\bigg(\|\bar\partial_t^\nu \nabla \rho_u^j\|_{L^2}+\|\bar\partial_t^\nu \rho_v^j\|_{L^2}\bigg)^2\\
&\leq CN\sum_{j=m+1}^r\lambda_j\sum_{j=1}^{n}\tau_j^{\alpha}\\
&\leq CnN^{1-\alpha}\sum_{j=m+1}^r\lambda_j.
\end{align*}
Then, we have
\begin{equation}\label{pod-5}
\frac{1}{N}\sum_{n=1}^N\|\theta_u^n\|_{L^2}^2
\leq  CN^{-\alpha}\sum_{j=m+1}^r\lambda_j\sum_{n=1}^Nn
\leq  CN^{2-\alpha}\sum_{j=m+1}^r\lambda_j.  
\end{equation}
Combining (\ref{pod-1}) and (\ref{pod-5}), the error between $e_m^n=\widetilde{U}_h^n-\widetilde{U}_m^n$ is given
\begin{align*}
\frac{1}{N}\sum_{n=1}^N\|\widetilde{U}_h^n-\widetilde{U}_m^n\|_{L^2}^2\leq  \frac{2}{N}\sum_{n=1}^N(\|\rho_u^n\|_{L^2}^2+\|\theta_u^n\|_{L^2}^2)\leq CN^{2-\alpha}\sum_{j=m+1}^r\lambda_j.    
\end{align*}
\end{proof}

\begin{thebibliography}{99}
\bibitem[LiaoH2]{LiaoH2018L1} H.~Liao, D.~Li, and J.~Zhang. \newblock Sharp error estimate of a nonuniform l1 formula for time-fractional reaction subdiffusion equations. \newblock {\em SIAM Journal on Numerical Analysis}, 56:1112--1133, 2018.
\bibitem[LyuP20]{LyuP2022SFOR} P.~Lyu and S.~Vong. \newblock A symmetric fractional-order reduction method for direct nonuniform approximations of semilinear diffusion-wave equations. \newblock {\em Journal of Scientific Computing}, 93:34, 2022.
\end{thebibliography}
\end{document}
